\section{数值验证}
\label{sec:simulation}

本节按第~\ref{sec:controller}~节的分段时序验证定理~\ref{thm:obs1}、\ref{thm:sta}、\ref{thm:ppc}、\ref{thm:pdf1}~及其合成（定理~\ref{thm:main}），以及命题~\ref{prop:stage2}、\ref{prop:approx}。全部结果由 \texttt{simulation/main\_ppc\_pdf\_verify.m} 生成（第~\ref{sec:sim_sweep}~节的 $T_o$ 扫描由 \texttt{simulation/main\_ppc\_to\_sweep.m} 生成），可复现。

\subsection{被控对象与仿真设置}
\label{sec:sim_model}

取 \(n=7\) 的自由漂浮空间机械臂，物理参数、初始状态、扰动力矩与关节力矩限幅取自文献~\cite{dou2025fractional}~的表~1、表~2 与式~(74)，以便与分数阶非奇异终端滑模（FONTSM）控制器在同一对象、同一工况下对比。仿真直接在 \(6+n=13\) 维浮基动力学上进行：每一采样时刻由完整惯性矩阵经 Schur 补得 \(M_e\) 与 \(C_e\)，控制器与被控对象之间无模型简化（该模型已与 Robotics System Toolbox 交叉校核，正运动学误差与 \(\dot H-2C\) 反对称性误差均达 \(10^{-16}\) 量级）。参考轨迹为关节空间五次多项式，\(T_{\mathrm{ref}}=8\,\mathrm s\)，幅值 \(10^{\circ}\)--\(20^{\circ}\)；采样步长 \(T_s=2\,\mathrm{ms}\)，仿真时长 \(T_f=20\,\mathrm s\)，\(\tau_{\max}=20\,\mathrm{N\cdot m}\)。初始关节误差为 \(0.2^{\circ}\)--\(0.5^{\circ}\)（\(3.5\times10^{-3}\)--\(8.7\times10^{-3}\,\mathrm{rad}\)）。

\textbf{算法实现。}观测器完全按式~\eqref{eq:ptdo} 在\emph{动量通道}上实现：由 \eqref{eq:momentum} 取 \(y=M_e(q_m)\dot q_m\)，由实际施加的力矩 \(\tau_{\rm act}\) 与解析模型构造
\begin{equation}
u_p=\tau_{\rm act}-C_e(q_m,\dot q_m)\dot q_m+\dot M_e(q_m,\dot q_m)\dot q_m,
\label{eq:sim_up}
\end{equation}
其中 \(\dot M_e\dot q_m\) 由 \(M_e\) 的中心差分的实质导数给出；估计量即扰动本身，故 \eqref{eq:computed_torque} 中的补偿量直接取 \(\hat d=z_2\) \eqref{eq:hatd}，无需再乘 \(M_e\)。每一拍按“先控制、后观测”的因果次序执行：

\begin{enumerate}
  \item[\textup{(i)}] 取已推进到 \(t_k\) 的估计 \(\hat d(t_k)=z_2\)，由 \eqref{eq:nu_schedule} 算 \(\nu(t_k)\)、由 \eqref{eq:computed_torque} 算 \(\tau_{\rm cmd}(t_k)\) 并限幅得 \(\tau_{\rm act}(t_k)\)；
  \item[\textup{(ii)}] 用 \(y(t_k)\)、\(\tau_{\rm act}(t_k)\) 与 \eqref{eq:sim_up} 把观测器 \eqref{eq:ptdo} 从 \(t_k\) 推进到 \(t_k+T_s\)。
\end{enumerate}
这样阶段 II 推进 \(\dot z_1=u_p+z_2+\kappa_1\sig^{1/2}(y-z_1)\) 所使用的零阶保持输入，正是区间 \([t_k,t_k+T_s)\) 内实际施加的力矩；若把 \textup{(i)(ii)} 的顺序颠倒（用上一拍的 \(\tau_{\rm act}\)），力矩在饱和边缘的逐拍跳变会作为未建模输入进入阶段 I 的奇异增益，并在 \(t\to T_o\) 处被放大成可观的估计偏差。

\textbf{采样实现的三个数值细节。}式~\eqref{eq:ptdo} 是连续时间命题，采样实现必须在以下三处作出选择；它们不改变算法，只决定命题~\ref{prop:approx}~中 \(\varepsilon_\Delta\) 与 \(\bar\Delta_r\) 的离散分量，故在此显式记录以便复现：

\begin{enumerate}
  \item[\textup{(a)}] \textbf{阶段 I 的压缩时间积分。}上行增益 \(6/(T_o-t)\)、\(6/(T_o-t)^2\) 在 \(t\to T_o\) 时奇异，显式 Euler 必然失稳。实现改在定理~\ref{thm:obs1}~证明所用的压缩时间坐标 \(\rho=\log\frac{T_o}{T_o-t}\) 下、以 \((z_1,z_2)\) 为状态积分
\[
\frac{\dd z_1}{\dd\rho}=-6z_1+(T_o-t)z_2+(T_o-t)u_p+6y,\qquad
\frac{\dd z_2}{\dd\rho}=\frac{6}{T_o-t}\,(y-z_1),
\]
用定步长 RK4（\(\Delta\rho\le0.05\)）推进，与 \(T_s\) 解耦；步内 \(y\) 与 \eqref{eq:sim_up} 的平滑部分用一阶预测 \(\dot y=u_p+\hat d\) 插值，使强迫项积分误差为 \(O(T_s^2)\)。
  \item[\textup{(b)}] \textbf{奇异增益正则化。}采样实现中 \(T_o-t\ll T_s\) 时，步内 \(y\) 的预测误差会被 \(1/(T_o-t)\) 放大。故在 \(T_o-t<s_{\min}=0.02\,\mathrm s\) 处把增益冻结为 \(6/s_{\min}\)、\(6/s_{\min}^2\)（即用 \(s_{\rm eff}=\max\{T_o-t,s_{\min}\}\)）。冻结后的误差系统仍线性稳定，其平衡残差为 \(e_2\approx\dot d\,s_{\min}=O(10^{-3})\,\mathrm{N\cdot m}\)，由阶段 II 的超螺旋在有限时间内收掉。
  \item[\textup{(c)}] \textbf{阶段 II 的离散切换。}下行用显式 Euler，并在每拍内细分为 \(n_{\rm sub}=5\) 个子步（\(\sig^{1/2}\) 与 \(\sgn\) 均取 Filippov 集内的单值，\(\sgn(0)=0\)）。细分不改变算法，只把 \(\sgn\) 的离散切换误差从 \(O(\kappa_2T_s)\) 降到 \(O(\kappa_2T_s/n_{\rm sub})\)。
\end{enumerate}

因此本节出现的“精确归零”应理解为连续时间极限；采样实现的稳态残差按命题~\ref{prop:approx}~报告，并在第~\ref{sec:sim_results}~节用 \(T_s\) 减半的实验核对其 \(O(T_s)\) 阶。

需要指出，该对象前四个关节的等效惯量为 \(1.2\times10^{2}\)--\(3.5\times10^{2}\,\mathrm{kg\cdot m^2}\)，而后三个关节仅 \(0.28\)--\(1.15\,\mathrm{kg\cdot m^2}\)，故 \(20\,\mathrm{N\cdot m}\) 限幅下关节 1--4 的加速度预算很窄，而关节 7 的残差 \(w_i\) 被 \(M_e^{-1}\) 放大约 \(3.6\) 倍。这两个事实分别决定了下文的饱和统计与关节 7 的贴边行为（第~\ref{sec:sim_infeasible}~节）。

\begin{figure}[t]
\centering
\begin{subfigure}{0.48\textwidth}
  \centering
  \includegraphics[width=\linewidth]{figures/ppc_base_config.pdf}
  \caption{惯性系下的构型}
  \label{fig:base_config}
\end{subfigure}\hfill
\begin{subfigure}{0.48\textwidth}
  \centering
  \includegraphics[width=\linewidth]{figures/ppc_base_position.pdf}
  \caption{基座质心位置}
  \label{fig:base_position}
\end{subfigure}
\caption{惯性系（初始基座系）中的运动（主算例 PPC，\(\rho_\infty=10^{-3}\)）：(a) 基座质心、各关节与末端在 \(t=0\)（实心）与 \(t=T_f\)（空心）的构型，以及基座质心与末端的轨迹；(b) 基座质心位置 \(r_b(t)\) 的三个分量。由零动量约束 \eqref{eq:gjm} 反算基座速度并积分得到（\texttt{simulation/make\_base\_motion\_figure.m}）。}
\label{fig:base_motion}
\end{figure}

图~\ref{fig:base_motion}(a) 给出惯性系下的初始/终止构型与运动轨迹，图~\ref{fig:base_motion}(b) 给出基座质心位置。整个任务中基座质心在惯性系中移动 \(100.6\,\mathrm{mm}\)（终点分量 \(51.4,\ 84.2,\ -19.8\,\mathrm{mm}\)，峰值速度 \(23.3\,\mathrm{mm/s}\)），姿态变化 \(0.193\,\mathrm{rad}\)，末端在惯性系中的行程为 \(1.05\,\mathrm{m}\)。相对约 \(3.5\,\mathrm{m}\) 的臂展，基座位移仅为 3\%，但它在惯性系中造成的构型偏移与 \eqref{eq:ejm} 的 Schur 补项 \(H_{bm}^{\top}H_b^{-1}H_{bm}\) 同源，不能按“基座固定”处理；这正是第~\ref{sec:model}~节采用等效关节空间模型并对 \(M_e\) 做 Schur 补的理由，也是第~\ref{sec:sim_infeasible}~节执行器预算与关节 7 贴边行为的物理背景。

\subsection{控制器与观测器参数}
\label{sec:sim_params}

\begin{table}[H]
\centering
\caption{控制器与观测器参数}
\label{tab:sim_params}
\small\begin{tabular}{p{3.2cm}p{5.8cm}p{4.2cm}}
\toprule
参数 & 取值 & 说明 \\
\midrule
预设时间 \(T_o\) & \(5\,\mathrm s\) & 总预设时间 \(T_o+2h=7\,\mathrm s\) \\
延迟 \(h\)，权重 \(R_h\) & \(1\,\mathrm s\)，\(\sin^4\!\bigl(\pi(\theta-h)/h\bigr)\) & \(C^{3}\) 启停 \\
性能函数 & 式~\eqref{eq:rho}（五次 smootherstep） & \(\rho_i\in C^{2}\) \\
包络端点 & \(\rho_{i0}=2\lvert e_i(0)\rvert\)，两组 \(\rho_{i\infty}\) & \(10^{-3}\) 与 \(2\times10^{-3}\) \\
漏斗增益 & \(k_{1i}=0.5\)，\(k_{2i}=1\) & \(K_p^{\mathrm{eff}}=0.5\)，\(K_d^{\mathrm{eff}}=1.5\) \\
基准增益 \(K_0\) & \([-K_p,\,-K_d]\) & \(A_c\) Hurwitz \\
\(K_p\)（关节 1--7） & 30, 30, 30, 25, 25, 20, 8 & 与既有仿真一致 \\
\(K_d\)（关节 1--7） & 12, 12, 12, 10, 10, 9, 4 & 同上 \\
观测器阶段 I & 式~\eqref{eq:ptdo} 上行，\(s_{\min}=0.02\,\mathrm s\) & 正则化奇异增益 \\
观测器阶段 II & 式~\eqref{eq:ptdo} 下行，\(\kappa_1=4\sqrt{L}\)，\(\kappa_2=2L\) & 满足 \eqref{eq:cond_gain} \\
交接整形 & \(\delta=\min\{0.4\,\mathrm s,0.15T_o\}\)，\(k_v=8/\delta\) & 第~\ref{sec:sim_mono}~节规则 R1 \\
离散化 & \(n_{\rm sub}=5\) 子步，\(\Delta\rho\le0.05\) & 见第~\ref{sec:sim_model}~节 \\
\bottomrule
\end{tabular}
\end{table}

观测器增益由假设~\ref{as:observer}~的 \textup{(C3)(C4)} 离线确定。由式~(74) 的解析导数在 \([0,T_f]\) 上取上界得
\[
L_i=\sup\lvert\dot d_i\rvert=[0.0185,\,0.0207,\,0.0201,\,0.0175,\,0.0100,\,0.0100,\,0.0100]\ \mathrm{N\cdot m/s},
\]
\(H_i=\sup\lvert\ddot d_i\rvert\le2.1\times10^{-3}\,\mathrm{N\cdot m/s^2}\)。按定理~\ref{thm:sta}~的注记取 \(\kappa_{2i}=2L_i\)、\(\kappa_{1i}=4\sqrt{L_i}\)，即 \(\kappa_{1i}\in[0.400,0.576]\)、\(\kappa_{2i}\in[0.0200,0.0415]\)，逐关节核验 \eqref{eq:cond_gain}：
\[
\kappa_{1i}^{2}-4L_i\frac{\kappa_{2i}+L_i}{\kappa_{2i}-L_i}
=[0.074,\,0.083,\,0.080,\,0.070,\,0.040,\,0.040,\,0.040]>0 ,
\]
故 \textup{(C4)} 全部满足（注意 \(L_i\) 只用于选增益，\eqref{eq:ptdo}~的在线实现并不需要它）。按推论~\ref{cor:exponential}~的整定准则 \(2\min\{k_{1i},k_{2i}\}T_o=5.0\) 大于 \(\ln\max_i(\rho_{i0}/\rho_{i\infty})=2.86\)（\(\rho_{i\infty}=10^{-3}\)），即 \(\rho\) 的收缩节奏与漏斗增益相容。

\(k_1,k_2\) 与基准增益 \(K_0\) 是两套独立参数：\(K_0\) 只在阶段 II 起作用，而 \(k_1,k_2\) 决定阶段 I 的物理等效 PD 增益（\(K_p^{\mathrm{eff}}=k_1k_2=0.5\)，\(K_d^{\mathrm{eff}}=k_1+k_2=1.5\)），因此阶段 I 的控制权柄主要由式~\eqref{eq:nu_ppc}~的\emph{障碍项}提供。这一结构特点在下面的结果讨论中会体现为“\(\tau_{\max}\) 与包络宽度的相容性”。

\subsection{离线判据}
\label{sec:sim_offline}

按第~\ref{sec:algorithm}~节第 (5) 步与命题~\ref{prop:stage2}~计算：\(\operatorname{cond}(W_c)=4.91\times10^{11}\)，单周期映射 \(\lVert\Delta_K(h)\rVert=3.00\times10^{-5}\)（相对 \(e^{A_c2h}\) 为 \(3.0\times10^{-5}\)）。尽管条件数很差，\texttt{pinv} 的截断正则化使零化残差仍很小，故 \(\varepsilon_\Delta=3.0\times10^{-5}\)、\(1/(1-\varepsilon_\Delta)\approx1.000\)。其余常数：\(\Gamma_h=0.397\)，\(\mathcal E_h=2.045\)，\(\int_h^{2h}\lVert K_c\rVert\,\mathrm ds=0.850\)，闭式上界 \(M_h^{\mathrm{ap}}=7.73\)，精确算子范数 \(M_{op}=2.045\)；逐关节位置行块 \(M_{h,i}=1.00\)（\(i=1,\dots,7\)），故命题~\ref{prop:stage2}~的判据退化为 \(\lVert z(T_o)\rVert_\infty\le\rho_{i\infty}\) 这一逐分量比较（见下文）。

\subsection{结果}
\label{sec:sim_results}

四个算例的验收指标列于表~\ref{tab:metrics}，其中 M1a/M1b 分别为观测阶段与 PDF 阶段的包络违反时长（\(\sum_i\int\mathbb 1\{\lvert e_i\rvert>\rho_i\}\mathrm dt\)，单位 \(\mathrm{s\cdot joint}\)），M2 为 \(T_o\) 时刻的进带深度 \(\max_i\lvert e_i(T_o)\rvert/\rho_{i\infty}\)，M4 为 PDF 驱动阶段的越界比 \(\max_{t\in[T_o,T_o+2h]}\max_i\lvert e_i\rvert/\rho_{i\infty}\)，M5 为 \(t\ge T_o+2h\) 后的 \(\sup\lVert z\rVert_\infty\)；末列为阶段 I 与全时段的力矩饱和占空比。

\begin{table}[H]
\centering
\caption{验收指标（\(T_o=5\,\mathrm s\)，\(h=1\,\mathrm s\)，\(\tau_{\max}=20\,\mathrm{N\cdot m}\)）。四个算例均含式~\eqref{eq:ptdo}~观测器与 PDF；PD 表示阶段 I 用基准反馈（不含第~\ref{sec:sim_mono}~节的整形项），其余用式~\eqref{eq:nu_ppc}~并含整形项。除标注 \(\dagger\) 者取 \(\rho_{i\infty}=10^{-3}\,\mathrm{rad}\) 外，均取 \(2\times10^{-3}\,\mathrm{rad}\)。}
\label{tab:metrics}
\small\setlength{\tabcolsep}{4pt}\begin{tabular}{p{5.0cm}cccccc}
\toprule
算例 & M1a & M1b & M2 & M4 & M5 & 饱和率(I/全) \\
\midrule
PD + PDF（消融 A）\(^{\dagger}\) & 9.80 & 0.44 & 4.38 & 4.38 & \(4.6\times10^{-5}\) & 22.8\%/5.7\% \\
PPC + PDF\(^{\dagger}\) & \textbf{0.178} & \textbf{0.000} & 0.907 & 0.907 & \(4.6\times10^{-5}\) & 1.72\%/0.43\% \\
PPC + PDF & 0.056 & \textbf{0.000} & 0.375 & 0.375 & \(4.6\times10^{-5}\) & 1.08\%/0.27\% \\
PPC + PDF（无扰动）\(^{\dagger}\) & \textbf{0.000} & \textbf{0.000} & 0.026 & 0.027 & \(2.9\times10^{-5}\) & 0\%/0\% \\
\bottomrule
\end{tabular}
\end{table}

\begin{figure}[t]
\centering
\includegraphics[width=\textwidth]{figures/ppc_envelope.pdf}
\caption{七关节跟踪误差 \(e_i(t)\) 与预设包络（红虚线 \(\pm\rho_i(t)\)、黑点线 \(\pm\rho_{i\infty}\)；竖实线 \(T_o=5\,\mathrm s\)，竖点划线 \(T_o+2h=7\,\mathrm s\)；\(\rho_\infty=10^{-3}\,\mathrm{rad}\)）。右下为误差包络总览。}
\label{fig:envelope}
\end{figure}

\begin{figure}[t]
\centering
\includegraphics[width=\textwidth]{figures/ppc_timeline.pdf}
\caption{三段时序指标：上为 \(\lVert z\rVert_\infty\)，中为归一化包络比 \(\max_i\lvert e_i\rvert/\rho_i\)（须 \(<1\)，红虚线为界），下为观测残差 \(\lVert w\rVert_\infty=\lVert M_e^{-1}(d-\hat d)\rVert_\infty\)（对数，阶段 I 的非匹配输入）。}
\label{fig:timeline}
\end{figure}

\begin{figure}[t]
\centering
\includegraphics[width=0.92\textwidth]{figures/ppc_phase1_compare.pdf}
\caption{阶段 I 对照：纯基准反馈 \(K_0z\)（蓝）与漏斗型 PPC（其余）。观测器峰值期间前者误差达 \(6.3\times10^{-2}\,\mathrm{rad}\)（初始误差的 7.2 倍），后者基本保持在 \(\rho_i(t)\) 内。}
\label{fig:phase1}
\end{figure}

\begin{figure}[t]
\centering
\includegraphics[width=0.95\textwidth]{figures/ppc_observer.pdf}
\caption{观测器（式~\eqref{eq:ptdo}，\(\rho_\infty=10^{-3}\) 算例）：三个关节的扰动真值 \(d_i\)（黑虚）与观测值 \(\hat d_i=z_{2i}\)（蓝），及逐拍最大观测误差（对数）。阶段 I 误差按压缩时间指数收缩、在 \(T_o\) 处达到 \(O(10^{-3})\)，阶段 II 的超螺旋把它压到 \(10^{-4}\) 量级并保持。}
\label{fig:observer}
\end{figure}

由表~\ref{tab:metrics}~与图~\ref{fig:envelope}--\ref{fig:observer}~可得以下结论。

\paragraph{(1) 阶段 I 的包络保持（定理~\ref{thm:funnel}、推论~\ref{cor:pd_compare}）}
PPC 的包络违反时长为 0.178\,s\(\cdot\)joint（占 \(5\,\mathrm s\times7\) 关节总时长的 \(0.5\%\)），归一化比峰值为 1.11；而同一包络下纯基准反馈违反 9.80\,s\(\cdot\)joint、归一化比峰值达 11.97（出现在 \(t\approx3.9\,\mathrm{s}\)）。这正是推论~\ref{cor:pd_compare}~的定量对照：给定同一条 \(\rho(t)\)，PD 的误差随观测残差线性恶化，而漏斗律把它压回包络内。进一步看，越界只发生在第 7 关节（0.172\,s\(\cdot\)joint）与第 6 关节（0.006\,s\(\cdot\)joint），其余关节为 0：该关节等效惯量仅 \(0.28\,\mathrm{kg\cdot m^2}\)，残差 \(w_7=(d_7-\hat d_7)/M_{77}\) 被放大约 3.6 倍，而 \(\rho_{70}=7\times10^{-3}\,\mathrm{rad}\)，故障碍项在 \(|e_7|\to\rho_7\) 附近的高增益被 \(2\,\mathrm{ms}\) 采样离散化，形成贴边的小幅极限环（图~\ref{fig:envelope}~右下）。这不是定理~\ref{thm:funnel}~的反例：该定理是连续时间结论且要求无饱和，而其贴合系数在此工况下退化（见下段），故采样实现只能保证“不破包络到离散化精度”。数值上该精度为归一化比的 \(11\%\)，出现在 0.5\% 的时段内。

\paragraph{(2) 预设时间进带（定理~\ref{thm:funnel}~(2)）}
\(T_o=5\,\mathrm s\) 时刻七个关节的 \(\lvert e_i(T_o)\rvert\) 为 \(2.1\times10^{-5}\)--\(9.1\times10^{-4}\,\mathrm{rad}\)（\(\rho_\infty=10^{-3}\)），即误差带深度的 \(2.1\%\!-\!90.7\%\)；\(\rho_\infty=2\times10^{-3}\) 时为带深的 \(0.13\%\!-\!37.5\%\)。需要说明，定理~\ref{thm:funnel}~给出的 \(\vartheta_i\) 在该工况下退化为 1：因为“oracle”残差上界取实测 \(\bar w=0.307\,\mathrm{rad/s^2}\)（由阶段 I 初始峰值决定），代入式~\eqref{eq:V_c}~得 \(c_i\sim10^{5}\)、\(\bar V_i\sim4.7\times10^{4}\)。故\emph{定理只保证不破包络、不保证深度}；实测深度远优于该（保守）界，与注~\ref{rem:unknown_w}~的解释一致。进带并不控制速度误差：第~\ref{sec:sim_mono}~节的交接整形把 \(\lVert\dot e(T_o)\rVert_\infty\) 由 \(3.2\times10^{-2}\) 压到 \(2.5\times10^{-4}\,\mathrm{rad/s}\)，使 PDF 阶段的位置误差可以直接单调下降（见结论 (4)）。

\paragraph{(3) 观测器：阶段 I 收敛与阶段 II 保持（定理~\ref{thm:obs1}、\ref{thm:sta}）}
图~\ref{fig:observer}~给出 \(d\) 与 \(\hat d=z_2\) 的对照。阶段 I 的残差 \(w\) 从 \(0.307\,\mathrm{rad/s^2}\) 按压缩时间指数下降（图~\ref{fig:timeline}~下图），在 \(T_o\) 处为 \(7\times10^{-4}\)；这与定理~\ref{thm:obs1}~(1) 的“收敛时刻恰为 \(T_o\)、精度由 \(L\) 与 \(T_o\) 决定”一致。阶段 II 的观测误差最大为 \(1.1\times10^{-2}\,\mathrm{N\cdot m}\)（切换后 \(0.1\,\mathrm{s}\) 内的瞬态），\(t\ge T_o+2h\) 后稳定在 \(2\times10^{-4}\)--\(4\times10^{-4}\,\mathrm{N\cdot m}\)，对应 \(w\) 的 \(10^{-4}\,\mathrm{rad/s^2}\) 量级。该残差并非定理~\ref{thm:sta}~的连续时间结论，而是第~\ref{sec:sim_model}~节 (c) 中 \(\sgn\) 的离散切换和 \(T_s\) 采样造成的；下一段关于 \(T_s\) 的实验给出其 $O(T_s)$ 阶的证据。

\paragraph{(4) PDF 阶段的单调收敛（命题~\ref{prop:stage2}、第~\ref{sec:sim_mono}~节规则）}
按第~\ref{sec:sim_mono}~节的规则整形交接状态后，PDF 阶段不再出现“先增后减”：M4\(=0.907<1\)、M1b\(=0\)，且 \(\lVert e(t)\rVert_\infty\) 在 \([T_o,T_o+2h]\) 上单调、各关节无超调（相对增量 \(8.9\times10^{-6}\rho_\infty\)；个别关节过零后回升 \(\le4.5\%\rho_\infty\)，见第~\ref{sec:sim_mono}~节 (4)）；\(\rho_\infty=2\times10^{-3}\) 时为 M4\(=0.375\)、M1b\(=0\)。作为对照，若去掉整形项（第~\ref{sec:sim_mono}~节的 R1），同一算例的 \(\dot e_7(T_o)=3.2\times10^{-2}\,\mathrm{rad/s}\) 会使 M4\(=5.38\)、M1b\(=1.18\)，与旧版结果一致：\emph{越界不是 PDF 结构的问题，而是交接速度的问题}。

消融 A（PD）不含整形项，其 PDF 阶段仍有 0.44\,s\(\cdot\)joint 的越界与 M4\(=4.38\)，可作为“无交接整形”的对照。命题~\ref{prop:stage2}~的判据仍按引理~\ref{lem:handover}~的 \(B_o=0.81\) 给出 813（因 \(\vartheta_i=1\) 退化），比实测保守约 900 倍；把实测 \(\lVert z(T_o)\rVert_2=1.17\times10^{-3}\) 代入则判据为 1.17，与实测 0.907 同阶——这说明判据的保守性来自 \(B_o\) 的退化，而非算子范数 \(M_{h,i}\)（实测 \(M_{h,i}=1.00\)）。整形后的交接状态使该判据几乎变紧，也从另一侧面说明 R1 的作用是把交接状态限制到定理允许的集合内。

\paragraph{(5) 归零与稳态残差（定理~\ref{thm:pdf1}、命题~\ref{prop:approx}）}
四个算例的 \(\lVert z(t)\rVert_\infty\) 均在 \(T_o+2h=7\,\mathrm{s}\) 前降到 \(4\times10^{-5}\) 以下（PDF 窗口内的最小值为 \(2.9\times10^{-5}\)），此后 \(\sup\lVert z\rVert_\infty=4.6\times10^{-5}\)（有扰动）与 \(2.9\times10^{-5}\)（无扰动），终端误差 \(3.8\times10^{-5}\,\mathrm{rad}\) 与 \(1.3\times10^{-5}\,\mathrm{rad}\)。为确认这是离散残差而非结构误差，固定其余参数把 \(T_s\) 由 \(2\,\mathrm{ms}\) 减半两次（\(\rho_\infty=2\times10^{-3}\)）。结果以三元组 \((T_s;\ \sup\lVert z\rVert_\infty;\ \lvert e(T_f)\rvert)\) 列为
\[
\begin{gathered}
(2\,\mathrm{ms};\ 4.6\times10^{-5};\ 2.8\times10^{-6}),\\
(1\,\mathrm{ms};\ 2.3\times10^{-5};\ 1.4\times10^{-6}),\\
(0.5\,\mathrm{ms};\ 1.2\times10^{-5};\ 4.5\times10^{-7}),
\end{gathered}
\]
两者均随 \(T_s\) 线性下降。这与命题~\ref{prop:approx}~的“\(\varepsilon_\Delta\) 与 \(\bar\Delta_r\) 决定最终界”一致，也说明本实现达到的是 \(O(T_s)\) 的最终有界性，而非定理~\ref{thm:main}~(iv) 的连续时间精确归零。无扰动算例的 M1a=M1b=0 且 M2\(=0.016\)，与 \(c_i=0\) 时推论~\ref{cor:exponential}~的指数进带图像吻合。

\paragraph{(6) 执行器负担}
等效物理增益较软的漏斗律（含交接整形）在同一工况下把阶段 I 的饱和占空比从 22.8\% 降到 1.1\%--1.7\%（全程 5.7\% 降到 0.3\%--0.4\%）。原因是阶段 I 的控制量由障碍项按需提供，误差未逼近包络时 \(\nu\) 很小，从而把有限的力矩预算留给参考轨迹的加速段；饱和的关节集中在惯量较大的 1--4（单关节 0.6\%--1.7\%），关节 5、6 各 0.12\%，第 7 关节全程不饱和。

\subsection{进入 PDF 阶段的单调收敛规则}
\label{sec:sim_mono}

本节把“单调收敛”精确定义为：\(\lVert e\rVert_\infty\) 在 \([T_o,T_o+2h]\) 上单调下降，且逐关节无超调 \(\lvert e_i(t)\rvert\le\lvert e_i(T_o)\rvert\)。点对点的 \(\lvert e_i\rvert\) 严格单调是更强的要求，在第~\ref{sec:pdf}~节的二阶误差通道与全耦合 Gramian 注入下不能由参数保证（见本节 (4)）。

若不施加交接整形，第 7 关节的误差在进入 PDF 阶段后会先由 \(4.7\times10^{-4}\,\mathrm{rad}\) 增至 \(5.4\times10^{-3}\,\mathrm{rad}\) 再归零。原因不在 PDF 增益：交接时刻 \(\dot e_7(T_o)=3.2\times10^{-2}\,\mathrm{rad/s}\) 与 \(e_7(T_o)\) 同号，故
\begin{equation}
\frac{\dd}{\dd t}\lvert e_i\rvert\Big|_{t=T_o}
=\operatorname{sgn}(e_i(T_o))\,\dot e_i(T_o)>0 ,
\label{eq:mono_necessary}
\end{equation}
即 \(\lvert e_7\rvert\) 在 \(T_o\) 处\emph{必然}上升，任何 \(K_0,K_c,h\) 都无法改变。由此得到下面的必要条件与三步参数规则。

\paragraph{(1) 必要条件（交接符号条件）}
\begin{equation}
e_i(T_o)\,\dot e_i(T_o)\le0,\qquad i=1,\dots,n .
\label{eq:handover_sign}
\end{equation}
更精确地，若 \(\dot e_i(T_o)=0\) 而 \(e_i(T_o)\neq0\)，则 \(e_i\ddot e_i(T_o)=-K_{p,i}e_i^2<0\)，\(\lvert e_i\rvert\) 从 \(T_o\) 起严格下降。

\paragraph{(2) 三步参数规则}
\textbf{R1（交接整形）：}在阶段 I 末段把辅助输入改为
\begin{equation}
\nu=\nu^{\mathrm{PPC}}-k_v\dot e,\qquad t\in[T_o-\delta,T_o),
\qquad \delta=\min\{0.4\,\mathrm s,\ 0.15T_o\},\quad k_v=8/\delta .
\label{eq:handover_damp}
\end{equation}
\(k_v\delta=8\gg1\) 使 \(\dot e_i(T_o)\to0\)（实测 \(\lVert\dot e(T_o)\rVert_\infty\le2.5\times10^{-4}\,\mathrm{rad/s}\)），从而满足 \eqref{eq:handover_sign}；\(\delta\le0.15T_o\) 保证整形项不压过漏斗本身的收缩。

\textbf{R2（自由传播窗口）：}\([T_o,T_o+h)\) 上 \(K_c\equiv0\)，每关节按 \(A_c\) 自由响应。由 \((e_i(T_o),0)\) 出发的位置响应无过冲的条件为 \(h\le\min_i t_i^*\)，其中
\begin{equation}
t_i^*=\begin{cases}
\dfrac{\pi/2+\arctan\!\bigl(\zeta_i/\sqrt{1-\zeta_i^2}\bigr)}{\sqrt{K_{p,i}}\,\sqrt{1-\zeta_i^2}}, & \zeta_i<1,\\[6pt]
\infty, & \zeta_i\ge1,
\end{cases}
\qquad
\zeta_i=\frac{K_{d,i}}{2\sqrt{K_{p,i}}} .
\label{eq:tstar}
\end{equation}
本文 \(K_0\) 给出 \(t_7^*=1.18\,\mathrm s\)（\(\zeta_7=0.707\)）、其余 \(t_i^*=\infty\)，故 \(h=1\,\mathrm s\) 满足并留 18\% 余量。

\textbf{R3（注入窗口校验）：}\([T_o+h,T_o+2h)\) 上的 Gramian 注入是全耦合的，且 \(\lVert\mathcal T(\theta,0)\rVert\) 可大于 1，故对任意交接方向不可能保证逐关节单调；须对实际 \(z(T_o)\) 做一次线性闭环校验（无需对象模型）：从测得的 \(z(T_o)\) 积分 \(\dot z=A_cz-BK_c(\theta)z(\theta-h)\)，要求
\begin{equation}
\max_{\theta\in[0,2h]}\lVert e(\theta)\rVert_\infty
\le(1+\varepsilon)\lVert e(T_o)\rVert_\infty,
\qquad \varepsilon=10^{-2},
\label{eq:mono_check}
\end{equation}
或逐关节增量 \(\le\varepsilon_\rho\rho_{i\infty}\)（\(\varepsilon_\rho=0.02\)）。不通过时增大 \(k_v\)（或 \(k_1,k_2\)）、放宽 \(\rho_{i\infty}\) 或减小 \(h\)。

\paragraph{(3) 效果}
按 R1--R3 更新后的主算例（表~\ref{tab:metrics}~第 2 行）M4\(=0.907\)、M1b\(=0\)，且 \eqref{eq:mono_check}~成立：\(\lVert e\rVert_\infty\) 在 \([T_o,T_o+2h]\) 上单调下降，各关节峰值不超过其交接值（以 \(\rho_{i\infty}\) 计的最大增量 \(8.9\times10^{-6}\)）；\(\rho_\infty=2\times10^{-3}\) 时 M4\(=0.375\)、最大增量 \(3.3\times10^{-3}\rho_\infty\)。阶段 I 的饱和占空比由 1.96\% 降至 1.72\%，包络违反时长由 0.21 降至 0.178\,s\(\cdot\)joint。

\paragraph{(4) 点对点单调的边界}
\eqref{eq:mono_check}~与表~\ref{tab:metrics}~保证的是范数与峰值包络，而不是每个 \(\lvert e_i\rvert\) 的点对点单调。原因有二：\textup{(i)} 二阶通道被驱向零时，若在过零时刻速度非零，则 \(\lvert e_i\rvert\) 在过零处取局部极小、随后回升；\textup{(ii)} Gramian 注入是跨关节耦合的（\(W_c^{-1}\) 满阵），初始为零的分量也会被其他关节激励。实测主算例中有 3 个关节（1、5、7）在 \(\theta\approx h\)（\(t\approx6.0\,\mathrm{s}\)）处过零，过零后 \(\lvert e_i\rvert\) 分别回升至 \(1.2\times10^{-5}\)、\(6.4\times10^{-6}\)、\(4.5\times10^{-5}\,\mathrm{rad}\)（最大为误差带的 4.5\%），随后再衰减到零。因此：\emph{在 \([T_o,T_o+h)\) 上逐点单调可证（R2）；在 \([T_o+h,T_o+2h)\) 上只能保证范数单调与无超调，点对点严格单调对任意交接方向都不成立}——这与命题~\ref{prop:stage2}~只给出 \(\lVert e_i\rVert\le M_{h,i}B_o\) 的结论一致，也是“精确归零”与“逐点单调”在二阶通道上的固有折中。

\paragraph{(5) 与定理的关系}
R1 是数值整形项，不属于定理~\ref{thm:main}~的连续时间闭环；它只把交接状态限制到定理允许的集合内（\(\bar w\) 与 \(B_o\) 的界仍成立），R3 则是与命题~\ref{prop:stage2}~同类的离线充分性校验。因此“精确归零”与“阶段 I 包络”的理论结论不变，规则只改善采样实现的暂态品质。

\subsection{不可行性与工程边界}
\label{sec:sim_infeasible}

\paragraph{(a) 包络必须与观测器残差和交接速度相容}
阶段 I 的残差 \(w\) 是 \(T_o\) 之前误差的主要来源：实测峰值 \(0.307\,\mathrm{rad/s^2}\)，其二次积分给出 \(\sim0.06\,\mathrm{rad}\) 量级的偏移（图~\ref{fig:phase1}），因此\emph{任何} \(\rho_0\lesssim10^{-2}\,\mathrm{rad}\) 的包络在该对象上都只能靠障碍项“硬压”。一个\emph{无补偿}的保守估计为
\begin{equation}
\rho_{i0}\ \gtrsim\ \tfrac12\lVert w\rVert_\infty\,T_o^{2},
\label{eq:feasible_envelope}
\end{equation}
它与第~\ref{sec:algorithm}~节关于 \(T_o\) 折中的讨论一致，并说明“预设包络”与“预设观测时间”不是两个可以独立选取的参数。式~\eqref{eq:feasible_envelope}~在该对象上过于保守（因观测器在收敛过程中已部分补偿），实际阈值由第~\ref{sec:sim_sweep}~节的扫描给出。另一端，交接速度 \(\lVert\dot e(T_o)\rVert\) 不被漏斗律本身约束：若不施加第~\ref{sec:sim_mono}~节的整形，它可达 \(3.2\times10^{-2}\,\mathrm{rad/s}\)，按 \(A_c\) 的惯性外扩约 \(5\times10^{-3}\,\mathrm{rad}\)，当 \(\rho_{i\infty}\lesssim5\times10^{-3}\,\mathrm{rad}\) 时命题~\ref{prop:stage2}~的判据会被破坏。R1 把该速度压到 \(2.5\times10^{-4}\,\mathrm{rad/s}\)，从而消除这一冲突；但整形窗口 \(\delta\le0.15T_o\) 与 \(k_v=8/\delta\) 的约束随 \(T_o\) 减小而变紧，这也是第~\ref{sec:sim_sweep}~节中 \(T_o=1\,\mathrm{s}\) 失效的另一原因。

\paragraph{(b) 初始包络过窄 + 力矩饱和 \texorpdfstring{$\Rightarrow$}{=>} 结构性失效}
定理~\ref{thm:funnel}~的包络保证只要求 \(\lvert\xi_i(0)\rvert<1\) 与残差有限，但在采样实现中“\(\xi_i(0)\) 接近 1”意味着障碍项在 \(t=0\) 附近就给出 \(10^{3}\) 量级的加速度指令，直接被 \(20\,\mathrm{N\cdot m}\) 限幅截断。实测三档初始包络：\(\rho_{i0}=2\lvert e_i(0)\rvert\)（\(\xi_i(0)=0.50\)）与 \(1.2\lvert e_i(0)\rvert\)（\(\xi_i(0)=0.83\)）均正常工作，阶段 I 饱和占空比分别为 \(1.7\%\) 与 \(1.4\%\)；而 \(\rho_{i0}=1.05\lvert e_i(0)\rvert\)（\(\xi_i(0)=0.95\)）使阶段 I 全时段饱和（\(100\%\)），闭环漂移到 \(10^{2}\,\mathrm{rad}\) 量级（实测 \(3.5\times10^{2}\,\mathrm{rad}\)）。可见可行边界在 \(\xi_i(0)\approx0.9\)，而不是定理所允许的 1。这不是数值问题而是假设~\ref{as:ideal}~的直接后果：\emph{一旦无饱和条件被破坏，定理~\ref{thm:funnel}~的证明链（尤其是式~\eqref{eq:Vdot}~中障碍项与交叉项的精确相消）即不成立}。工程上可行的补救是把障碍项限幅（例如限在 \(\lvert\xi_i\rvert=0.95\) 处的值），使破包络后律退化为有界 PD 型反馈，避免灾难性漂移；代价是包络的\emph{严格}性变为“实际保包络”。

\paragraph{(c) Gramian 病态的现实影响}
\(\operatorname{cond}(W_c)=4.9\times10^{11}\) 虽未破坏 \(\Delta_K(h)\) 的数值零化（\(\lVert\Delta_K(h)\rVert=3.0\times10^{-5}\)），但使 \(K_c\) 的幅值对参数敏感：\(\int_h^{2h}\lVert K_c\rVert\,\mathrm ds=0.85\) 与 \(\mathcal E_h=2.05\) 共同决定 \(M_h^{\mathrm{ap}}=7.7\)，而后者是 M4 判据保守性的来源之一。若把 \(h\) 减小到 \(0.3\,\mathrm s\)，\(\operatorname{cond}(W_c)\) 升至 \(10^{13}\) 以上、\(\lVert\Delta_K(h)\rVert\) 超过 \(10^{-2}\)，PDF 的时域优势即消失（与既有结论一致）。故 \(h=1\,\mathrm s\) 是“零化精度”与“周期内放大”之间的折中。

\paragraph{(d) 奇异增益只能正则化实现}
式~\eqref{eq:ptdo} 上行的增益在 \(T_o\) 处发散，任何采样实现都无法把误差积到严格零；第~\ref{sec:sim_model}~节 (b) 的正则化 \(s_{\min}\) 是必要的，它把“预设时刻精确收敛”降级为“预设时刻 \(O(s_{\min})\) 收敛”。这一降级与命题~\ref{prop:approx}~的框架相容（残差进入 \(\bar\Delta_r\)），但意味着\emph{实现层面的收敛时刻精度由 \(s_{\min}\) 与 \(T_s\) 的比值决定}：\(s_{\min}/T_s\) 越小，交接残差越小而步内放大越强，二者不可同时最优。本文取 \(s_{\min}/T_s=10\) 作为折中。

\subsection{观测时间 \(T_o\) 的扫描}
\label{sec:sim_sweep}

与既有版本不同，本算法的观测器没有独立的内部时间常数 \(T_c\)：阶段 I 的收敛时刻就是 \(T_o\)，阶段 II 的增益只由 \(L_i\) 从下方约束。因此可整定的“时间”参数只剩 \(T_o\)。固定其余参数、令 \(T_o\in\{1,1.5,2,3,4,5,6,8\}\,\mathrm{s}\)，对 PPC 与纯 PD 两种阶段 I 各运行一次（共 16 次，脚本 \texttt{main\_ppc\_to\_sweep.m}）。为使扫描反映 \(T_o\) 本身而非包络可行性，取与观测瞬态相容的包络 \(\rho_{i0}=2\lvert e_i(0)\rvert\)、\(\rho_{i\infty}=2\times10^{-3}\,\mathrm{rad}\)。

\begin{table}[H]
\centering
\caption{\(T_o\) 扫描的关键点（\(\rho_\infty=2\times10^{-3}\,\mathrm{rad}\)，\(h=1\,\mathrm s\)）}
\label{tab:sweep}
\small\begin{tabular}{ccccccc}
\toprule
 & \multicolumn{3}{c}{PPC \(+\) PTDO \(+\) PDF} & \multicolumn{3}{c}{PD \(+\) PTDO \(+\) PDF} \\
\cmidrule(lr){2-4}\cmidrule(lr){5-7}
\(T_o\,[\mathrm s]\) & 阶段 I 饱和 & \(\sup\lVert z\rVert_\infty\) & M1a / M2 / M4 & 阶段 I 饱和 & \(\sup\lVert z\rVert_\infty\) & M1a \\
\midrule
1.0 & 53.6\% & \(1.4\times10^{0}\) & 1.05 / 41.0 / 571 & 87.6\% & \(4.6\times10^{-5}\) & 3.00 \\
1.5 & 29.6\% & \(5.0\times10^{-5}\) & 0.016 / 0.86 / 0.86 & 76.4\% & \(4.7\times10^{-5}\) & 5.48 \\
2.0 & 5.9\% & \(4.7\times10^{-5}\) & 0.064 / 0.85 / 0.85 & 57.2\% & \(4.7\times10^{-5}\) & 7.12 \\
3.0 & 2.4\% & \(4.5\times10^{-5}\) & 0.074 / 0.83 / 0.83 & 38.0\% & \(4.6\times10^{-5}\) & 7.14 \\
4.0 & 1.9\% & \(4.4\times10^{-5}\) & 0.070 / 0.90 / 0.90 & 28.5\% & \(4.4\times10^{-5}\) & 7.74 \\
5.0 & 1.1\% & \(4.6\times10^{-5}\) & 0.056 / 0.38 / 0.38 & 22.8\% & \(4.6\times10^{-5}\) & 8.03 \\
6.0 & 3.4\% & \(4.5\times10^{-5}\) & 0.21 / 0.87 / 0.87 & 19.0\% & \(4.3\times10^{-5}\) & 8.89 \\
8.0 & 1.0\% & \(4.5\times10^{-5}\) & 0.096 / 0.78 / 0.78 & 14.2\% & \(4.5\times10^{-5}\) & 10.8 \\
\bottomrule
\end{tabular}
\end{table}

\begin{figure}[t]
\centering
\includegraphics[width=\textwidth]{figures/ppc_to_sweep.pdf}
\caption{\(T_o\) 扫描。(a) 阶段 I 观测残差峰值（左轴）与阶段 I 饱和占空比（右轴）；(b) 稳态残差与终端误差；(c) 阶段 I 包络违反时长 M1a（左轴）与进带深度 M2、PDF 越界比 M4（右轴，PPC）；按第~\ref{sec:sim_mono}~节规则整形后 M4$\equiv$M2（不再超调）。}
\label{fig:sweep}
\end{figure}

由图~\ref{fig:sweep}~与表~\ref{tab:sweep}~可得：

\paragraph{(i) \(T_o\) 存在下界：包络收缩必须可由执行器实现}
\(T_o=1\,\mathrm s\) 时性能函数收缩过快（\(\lvert\dot\rho_i\rvert\propto1/T_o\)），阶段 I 饱和占空比达 53.6\%，漏斗被饱和击穿（M1a\(=1.05\)、M2\(=41\)），阶段 II 也无法把已失稳的状态拉回（\(\sup\lVert z\rVert_\infty=1.4\,\mathrm{rad}\)）。这与式~\eqref{eq:tc_tradeoff}~的“增益下界—采样/限幅上界”冲突同源：此处受约束的是 \(\rho\) 的收缩速率而非观测器增益；同时 R1 的窗口 \(\delta\le0.15T_o\) 与 \(k_v=8/\delta\) 在 \(T_o=1\,\mathrm s\) 处只能给出 \(\delta=0.15\,\mathrm s\)、\(k_v=53\)，整形项本身也已贴近限幅。\(T_o=1.5\,\mathrm s\) 时虽能收敛（M2\(=0.86\)），但阶段 I 饱和仍达 29.6\%，属于“可行但昂贵”的边界。

\paragraph{(ii) \(T_o\ge2\,\mathrm s\) 后稳态残差与 \(T_o\) 无关}
\(T_o\in[2,8]\,\mathrm s\) 的稳态残差都在 \(4.4\)--\(4.7\times10^{-5}\)，终端误差同为 \(3.8\)--\(4.2\times10^{-5}\)，即 \(O(T_s)\) 的离散地板而非 \(T_o\) 决定的误差；\(T_o\) 的增大只单调降低阶段 I 的饱和占空比（53.6\%\(\to\)1.0\%）与力矩裕度压力。另一方面总预设时间 \(T_{\mathrm{pre}}=T_o+2h\) 随 \(T_o\) 增大，故 \(T_o\in[2,5]\,\mathrm s\) 是该对象上的合理窗口。

\paragraph{(iii) PPC 相对 PD 的优势在每个 \(T_o\) 上都成立，且随 \(T_o\) 增大而扩大}
同一 \(T_o\) 下 PPC 的 M1a 比 PD 小 1.5--2 个数量级（\(T_o=5\) 时 0.056 对 8.03），阶段 I 饱和占空比小 15--25 倍。PD 的 M1a 随 \(T_o\) 单调增大（3.0\(\to\)10.8），因为它没有包络，误差峰值由残差作用时长决定；而 PPC 的 M1a 始终在 0.016--0.21 的小量级，不随 \(T_o\) 单调变化。这再次印证推论~\ref{cor:pd_compare}：漏斗律把峰值从“破坏精度”降级为“消耗进带余量”。

\paragraph{(iv) 在单调规则下 M4\(\equiv\)M2，整定目标归结为进带深度}
按第~\ref{sec:sim_mono}~节规则整形后，所有可行 \(T_o\) 的 PDF 阶段都不再超调，M4 与 M2 相等（表中三列 M1a/M2/M4 的后两列相同）。此时 M2 在 0.38--0.90 间非单调变化（\(T_o=5\) 时最小 0.38，\(T_o=4\) 时 0.90），反映的是观测残差 \(w(t)\) 与 \(\rho(t)\) 收缩节奏的相对相位，不能用单一趋势概括。工程上应以“M2 留足余量、阶段 I 饱和可接受”为准逐案校验，而不能只按 \(T_o\) 的大小排序。
